\subsection{Sample representativeness}
\label{sec:p1_representativeness}

\Cref{tab:2A:representativeness} and \cref{fig:2A:representativeness_region} compare the
interviewed and visited schools (300 and 102 of them matched to the school panel; 303 and
103 in all) to the national population of 3,312 open secondary and all-through schools in
England, across a range of structural characteristics drawn from the school panel spine and
the Get Information About Schools register, together with the inspection grade as at August
2024 and the pre-pandemic Progress~8 of 2018--19.

The interviewed and visited samples are broadly representative of the national population
on most structural dimensions. School type (academy, maintained, free school), phase of
education, gender composition, admissions policy, and urban--rural classification are all
closely matched to national proportions, with no statistically significant differences.
Mean school size and free school meals percentage are also similar to national averages:
interviewed schools average 1,109 pupils (national mean: 1,079) and a free school meals
share of 27.9\% (national mean: 29.9\%).

The main exception is regional distribution
(\cref{fig:2A:representativeness_region}). Both samples over-represent London and, for
interviews, the South East: taken together, London and the South East account for
44.7\% of interviewed schools versus 30.3\% nationally. The visited sample is
particularly concentrated geographically---London accounts for 45.1\% of visits compared
to 15.0\% of schools nationally---reflecting the logistical constraints of conducting
full-day school visits. Northern England is correspondingly under-represented: the North
West (13.9\% nationally, 8.7\% of interviewed, 2.9\% of visited), the West Midlands
(11.9\%, 7.7\%, 1.0\%), and Yorkshire and the Humber (9.6\%, 5.3\%, 0.0\%) all feature
substantially less in the sample than in the population. Visited schools are also more
likely to have a sixth form (nan\% versus nan\% nationally), consistent with the London
concentration given that London schools disproportionately offer post-16 provision.

Both samples also lean toward stronger schools: 25 per cent of visited schools held an
Outstanding grade as at August 2024 against 15 per cent nationally, and their mean
2018--19 Progress~8 was $+$0.17 against $+$0.04 nationally; the interviewed sample shows the
same lean more weakly (19 per cent and $+$0.08).

These patterns should be borne in mind when interpreting findings, particularly for the
visit-based tier. Warmth and strictness estimates for the primary tier
(\Cref{sec:p1_data}) may not be fully generalisable to schools in the North of England,
and cross-tier comparisons should account for the possibility that regional compositional
differences contribute to attenuation.
